\documentclass[a4paper]{article}

\usepackage{graphicx}
\usepackage{amsmath,amssymb}
\usepackage{minted}
\usepackage{multirow}
\usepackage{float}
\usepackage{todonotes}
\usepackage{forest}
\usepackage{xstring}
\usepackage{xcolor}
\usepackage[most]{tcolorbox}
\usepackage{xparse}
\usepackage{natbib}

\usetikzlibrary{graphs,graphdrawing}
\usegdlibrary{layered}

% for external graphviz
\input{/home/donkarlo/Dropbox/repo/nd_ai_project/src/nd_ai/knoweldge_representation/language/formal/computer/programming/tex/latex/code_clips/external_graphviz.tex}

% for tags
\input{/home/donkarlo/Dropbox/repo/nd_ai_project/src/nd_ai/knoweldge_representation/language/formal/computer/programming/tex/latex/parts/control_sequence/kind/semtag/semtag.tex}

% for links
\input{/home/donkarlo/Dropbox/repo/nd_ai_project/src/nd_ai/knoweldge_representation/language/formal/computer/programming/tex/latex/code_clips/seehere.tex}

\title{Guassian process regression}
\date{}
\author{Mohammad Rahmani}

\begin{document}
    \maketitle
    \cite{barfoot-2014-batch-continuous-time-trajectory-estimation-as-exactly-sparse-gaussian-process-regression} propose a continuous-time approach for batch trajectory estimation using Gaussian Process regression. The trajectory is treated as a one-dimensional Gaussian Process where time is the independent variable. This is useful for asynchronous or multi-rate robotic sensors, such as scanning laser rangefinders, rolling-shutter cameras, and other sensors that produce observations at different timestamps. Instead of using direct linear or spline interpolation between discrete poses, the method represents the whole trajectory as a continuous-time probabilistic function. The key contribution is that, for Gaussian Process priors generated by linear time-varying stochastic differential equations, the inverse kernel matrix becomes exactly sparse and block-tridiagonal. This allows efficient GP regression and interpolation, making it possible to query the robot state at any time of interest with low computational cost.
    \bibliographystyle{plainnat}
    \bibliography{/home/donkarlo/Dropbox/repo/nd_shared_research/bibliography/bibliography.bib}
\end{document}